\section{Results}
\label{sec:results}

\subsection{The skill effect is sparse on successful contexts}
Figure~\ref{fig:concentration} answers RQ1.  The top 1\% of positions capture 42.11\% of aggregate positive gain; the top 5\% capture 84.05\%; and the top 10\% capture 96.12\%.  All 61 trajectories individually satisfy $\mathcal{C}(10\%)\geq 0.6$ and $\mathcal{C}(20\%)\geq 0.8$.  Random selection closely tracks its nominal budget.  This strongly supports localization, although it does not establish controllability by a short prefix.

\begin{figure}[t]
\centering
\begin{tikzpicture}
\begin{axis}[
    width=0.98\linewidth,
    height=4.7cm,
    xmin=0,xmax=20,
    ymin=0,ymax=100,
    xtick={1,2,5,10,20},
    ytick={0,20,40,60,80,100},
    xlabel={selected token budget (\%)},
    ylabel={positive-gain mass (\%)},
    grid=major,
    legend style={font=\scriptsize,at={(0.98,0.05)},anchor=south east},
    tick label style={font=\scriptsize},
    label style={font=\small}
]
\addplot[color=skillblue,very thick,mark=*] coordinates {(1,42.11) (2,59.72) (5,84.05) (10,96.12) (20,99.76)};
\addlegendentry{Positive locator}
\addplot[color=skillgray,dashed,thick,mark=square*] coordinates {(1,0.99) (2,1.99) (5,4.99) (10,9.97) (20,19.84)};
\addlegendentry{Random}
\end{axis}
\end{tikzpicture}
\caption{Aggregate concentration on 61 successful training trajectories.  Skill-induced positive gold-token gain is highly concentrated under teacher forcing.}
\label{fig:concentration}
\end{figure}


Positive and Combined select overlapping but non-identical cores (mean Jaccard 0.697).  At 5\%, Positive captures more realized gain (81.91\% vs.\ 76.23\%), whereas Combined captures more JS mass (50.61\% vs.\ 44.69\%) and combined-score mass (97.56\% vs.\ 95.90\%).  Full-distribution information therefore shifts capacity toward positions where the skill changes alternatives beyond the observed token.

\subsection{Selective compression improves over full-trajectory SoftSkill}

\begin{table}[t]
\caption{Matched SpreadsheetBench test results.  All methods use the same frozen Qwen checkpoint, 280 tasks, greedy decoding, batch size 2, and an 8,192-token cap.  $\Delta$ is relative to full-trajectory SoftSkill.}
\label{tab:test-main}
\centering
\small
\begin{tabular}{lrrr}
\toprule
Method & Solved & Success (\%) & $\Delta$ (pp) \\
\midrule
No Skill & 97/280 & 34.64 & +4.29 \\
Full-trajectory SoftSkill & 85/280 & 30.36 & -- \\
PRCB-v5 & 96/280 & 34.29 & +3.93 \\
\textbf{Combined-5\%+Preserve} & \textbf{101/280} & \textbf{36.07} & \textbf{+5.71} \\
\bottomrule
\end{tabular}
\end{table}


Table~\ref{tab:test-main} answers RQ2.  Combined-5\%+Preserve solves 101/280 test tasks, 16 more than full-trajectory SoftSkill.  The paired difference is significant under an exact McNemar test ($p=.0479$).  The improvement is entirely on cell-level tasks (61 vs.\ 45 of 193); both solve 40 of 87 sheet-level tasks.  This suggests that sparse instruction transfer helps local edits more than global workbook transformations.

The strongest conclusion is relative to the matched full-trajectory baseline, not the base model.  No Skill solves 97 tasks, and the four-task net advantage of Combined is not significant ($40$ Combined-only vs.\ $36$ no-skill-only, $p=.731$).  PRCB-v5 reaches 96/280 and likewise does not outperform No Skill.  Selectivity recovers the degradation introduced by indiscriminate trajectory fitting, but does not yet robustly add skill beyond the base model's existing competence.

\subsection{Validation ablations}

\begin{table*}[t]
\caption{SpreadsheetBench validation ablations (40 tasks).  ``Shared'' means Positive and Combined use the identical matched preservation indices.  Runs marked incomplete are excluded.}
\label{tab:validation}
\centering
\small
\begin{tabular}{lrrrrl}
\toprule
Training target & CE tokens & Coverage & Preserve & Epochs & Solved \\
\midrule
Full trajectory & 54,990 & 100.00\% & -- & 3 & 11/40 (27.5\%) \\
Random-5\% & 2,838 & 5.06\% & -- & 3 & 15/40 (37.5\%) \\
Random-5\% + KL & 2,838 & 5.06\% & 2,777 & 1 & 15/40 (37.5\%) \\
Positive-5\% + KL & 2,838 & 5.06\% & 2,777 & 1 & \textbf{18/40 (45.0\%)} \\
Positive-5\% + shared KL & 2,838 & 5.06\% & 2,777 & 1 & 10/40 (25.0\%) \\
Combined-5\% + shared KL & 2,838 & 5.06\% & 2,777 & 1 & 16/40 (40.0\%) \\
Positive window L2/R8 + KL & 19,424 & 35.25\% & 2,777 & 1 & 12/40 (30.0\%) \\
\bottomrule
\end{tabular}
\end{table*}


Table~\ref{tab:validation} answers RQ3.  With independent preservation sets, Positive-5\%+Preserve reaches 18/40 versus 15/40 for matched Random-5\%+Preserve.  In the stricter paired comparison with the same preservation indices, Combined improves over Positive (16/40 vs.\ 10/40), but the exact paired test is inconclusive ($p=.146$).  Expanding each Positive core by a left-2/right-8 window increases CE coverage from 5.06\% to 35.25\% yet reduces success to 12/40.  More local context is therefore not automatically better; it reintroduces abundant low-specificity targets and changes the loss scale.

\subsection{Task-specific capacity oracle}

To separate shared-prefix interference from representational limits, we train one prefix for each of the 61 successful training tasks and freely regenerate the \emph{same} task.  Combined core-only selection at 10\% solves 35/61 (57.38\%), compared with 32/61 (52.46\%) for the hard textual skill and 26/61 (42.62\%) for No Skill.  Five percent solves 28/61 and 20\% solves 33/61.  This is an oracle memorization/closure test, not held-out evidence, but it shows that an eight-token prefix can sometimes turn teacher-forced fitting into a successful rollout when cross-task interference is removed.  The non-monotonic 5/10/20\% curve also rules out ``insufficient token coverage'' as the sole explanation.
